\subsection{Reinforcement Learning}\label{reinforcement-learning}

RL is a branch of machine learning characterized by two core components
(\colorbox{yellow!40}{Reference}): (i) the agent; and (ii) the environment. The agent
observes the current state of the environment, selects an action, and
receives feedback in the form of a reward associated with that action
(Sutton and Barto, 2018). Through iterative interactions with the
environment, the agent refines its decision-making strategy, or policy,
to maximize long-term rewards (Esteso \emph{et al.}, 2023). Unlike
traditional machine learning approaches that rely on static,
pre-existing datasets, RL agents learn from dynamic, continuously
evolving datasets generated through their ongoing interaction with a
virtual environment (Panzer and Bender, 2022). This environment is often
modeled using discrete-event simulation, which effectively captures
system dynamics and variability (Kuhnle \emph{et al.}, 2022; Nguyen
\emph{et al.}, 2022).

RL problems are typically framed as a sequence of interactions between
the agent and the environment at discrete time steps
\(t \in \left\{ 0,\ 1,\ 2,\ \ldots,T \right\}\) (\colorbox{yellow!40}{Reference}). At
each time step \(t\), the agent receives a representation of the
environment's state, \(S_{t}\mathcal{\in S}\), where \(\mathcal{S}\) is
the set of possible states (\colorbox{yellow!40}{Reference}). Based on this state, the
agent selects an action, \(A_{t}\mathcal{\in A(}S_{t})\), where
\(\mathcal{A(}S_{t})\) represents the set of actions available in state
\(S_{t}\) (\colorbox{yellow!40}{Reference}). One time step later, as a result of its
action, the agent receives a numerical reward,
\(R_{t + 1}\mathcal{\in R \subset}\mathbb{R}\), and transitions to a new
state, \(S_{t + 1}\) (\colorbox{yellow!40}{Reference}). In manufacturing applications of
RL, the environment corresponds to the manufacturing system itself (or a
model of it), such as an assembly line or warehouse (\colorbox{yellow!40}{e.g.,
Reference}). The agent acts as the decision-maker, which could be a
production scheduler, for example (\colorbox{yellow!40}{e.g., Reference}). The state
represents the current status of the system, such as the number of
machines that are idle, running, or under maintenance (\colorbox{yellow!40}{e.g.,
Reference}). The action refers to a decision the agent makes within the
system, such as assigning a new work order to a machine or adjusting the
priority of a pending job (\colorbox{yellow!40}{e.g., Reference}). The reward represents
the payoff the agent receives in response to its actions, typically
aimed at improving a specific performance measure, such as production
volume or lead time (\colorbox{yellow!40}{e.g., Reference}).

The environment in RL can be deterministic or stochastic. In the latter
case, it is often modeled as a finite Markov Decision Process (MDP),
where the environment's response at \(t + 1\) depends only on the state
and action representations at \(t\) (\colorbox{yellow!40}{Reference}). This assumption
simplifies the dynamics of agent-environment interactions, which can be
described by a distribution function, such as (\colorbox{yellow!40}{Reference}):

\begin{equation}
p\left( s',r \middle| s,a \right) = \Pr\left\{ S_{t + 1} = s',\ R_{t + 1} = r \middle| S_{t} = s,A_{t} = a \right\}
\tag{1}
\end{equation}

for all possible states \(s\), \(s'\), rewards \(r\), and actions \(a\).
This formulation allows prediction of the next state and expected reward
based on the current state and action, though it requires prior
knowledge on the distribution function (\colorbox{yellow!40}{Reference}).

The agent's behavior is governed by a policy,
\(\pi_{t}\left( a \middle| s \right)\), which represents the probability
of taking action \(A_{t} = a\) in state \(S_{t} = s\) (\colorbox{yellow!40}{Reference}).
If the policy involves no randomness, then only one action can be taken
from each state, resulting in a deterministic policy \(\pi_{t}(s)\)
(\colorbox{yellow!40}{Reference}). RL methods define how the agent updates its policy
based on its experiences, with the ultimate goal of maximizing the total
rewards received over time, known as the return \(G_{t}\):

\begin{equation}
G_{t} = \sum_{k = 0}^{T - t - 1}{\gamma^{k}R_{t + k + 1},}
\tag{2}
\end{equation}

where \(\gamma\) is the discount rate \((0 \leq \gamma \leq 1)\) that
determines the present value of future rewards (\colorbox{yellow!40}{Reference}). A
reward received \(k\) time steps in the future is worth only
\(\gamma^{k - 1}\) times its immediate value.

Most RL algorithms rely on estimating value functions, which assess the
`goodness' of being in a given state or taking a specific action
(\colorbox{yellow!40}{Reference}). This goodness is defined in terms of the expected
return \(\mathbb{E}_{\pi}\lbrack.\rbrack\) and is essential for
identifying optimal policies \(\pi_{*}\), where
\(\upsilon_{*}(s) \geq \upsilon_{\pi}(s)\) or
\(q_{*}(s,a) \geq q_{\pi}(s,a)\) for all \(s \in \mathcal{S}\) and
\(a \in \mathcal{A}\), (\colorbox{yellow!40}{cf., Reference}). The two primary value
functions for finite MDPs are (\colorbox{yellow!40}{Reference}):

\begin{enumerate}
\def\labelenumi{\roman{enumi}.}
\item
  State-value function \((\upsilon_{\pi}(s))\): The expected return
  starting in state \(s\) and following policy \(\pi\):
\end{enumerate}

\begin{equation}
\upsilon_{\pi}(s) = \mathbb{E}_{\pi}\left\lbrack G_{t}|S_{t} = s \right\rbrack = \mathbb{E}_{\pi}\left\lbrack \left. \ \sum_{k = 0}^{T - t - 1}{\gamma^{k}R_{t + k + 1}} \right|S_{t} = s \right\rbrack;
\tag{3}
\end{equation}

\begin{enumerate}
\def\labelenumi{\roman{enumi}.}
\setcounter{enumi}{1}
\item
  Action-value function \((q_{\pi}(s,a))\): The expected return starting
  from state \(s\), taking action \(a\), and subsequently following
  policy \(\pi\):
\end{enumerate}

\begin{equation}
q_{\pi}(s,a) = \mathbb{E}_{\pi}\left\lbrack G_{t}|S_{t} = s,A_{t} = a\  \right\rbrack = \mathbb{E}_{\pi}\left\lbrack \left. \ \sum_{k = 0}^{T - t - 1}{\gamma^{k}R_{t + k + 1}} \right|S_{t} = s,\ A_{t} = a \right\rbrack.
\tag{4}
\end{equation}

Value functions satisfy recursive relationships, which allows the use of
Bellman optimality equation (\colorbox{yellow!40}{cf., Reference}). This leads to the
optimal forms of the state-value function \({(\upsilon}_{*}(s))\) and
the action-value function \({(q}_{*}(s,a)\)), defined as follows
(\colorbox{yellow!40}{Reference}):

\begin{equation}
\upsilon_{*}(s) = \max_{a \in A(s)}{\sum_{s',r}^{}{p\left( s',r \middle| s,a \right)\left\lbrack r + \gamma\upsilon_{*}(s') \right\rbrack}.}
\tag{5}
\end{equation}

\begin{equation}
q_{*}(s,a) = \sum_{s',r}^{}{p\left( s',r \middle| s,a \right)}\left\lbrack r + \gamma\max_{a'}{q_{*}(s',a')} \right\rbrack.
\tag{6}
\end{equation}

The computation of optimal policies depends on how the agent predicts
the environment\textquotesingle s response to its actions
(\colorbox{yellow!40}{Reference}). In model-based methods, where the agent uses a model
to plan actions in advance, such as estimating
\(p\left( s',r \middle| s,a \right)\), dynamic programming algorithms
are typically employed (\colorbox{yellow!40}{e.g., Reference}). In contrast, model-free
methods focus on learning the direct relationship between actions and
returns through `trial and error', often utilizing temporal-difference
algorithms (\colorbox{yellow!40}{Reference}). Model-free methods can be further divided
into on-policy and off-policy strategies (\colorbox{yellow!40}{Reference}). An on-policy
strategy learns from the actions taken by the current policy, with SARSA
being the prominent learning algorithm (\colorbox{yellow!40}{e.g., Reference}), while an
off-policy strategy learns from actions taken by a combination of the
current and different policies, typically implemented using the
Q-learning algorithm (\colorbox{yellow!40}{e.g., Reference}). Furthermore, model-free
methods can be categorized into gradient-based and gradient-free
methods. Gradient-based methods, such as Proximal Policy Optimization
(PPO), optimize policies by directly adjusting parameters in the
direction of the gradient of the expected return with respect to the
policy parameters (\colorbox{yellow!40}{Reference}). The PPO algorithm specifically
improves upon traditional policy gradient methods by employing a clipped
objective function to prevent large updates that could destabilize
training (\colorbox{yellow!40}{Reference}). The core PPO objective is expressed as
follows:

\begin{equation}
L^{CLIP}(\theta) = \mathbb{E}_{t}\left\lbrack \min\left( p_{t}(\theta){\widehat{A}}_{t},\ clip(p_{t}(\theta),1 - \epsilon,1 + \epsilon){\widehat{A}}_{t} \right) \right\rbrack;
\tag{7}
\end{equation}

where
\(p_{t}(\theta) = \frac{\pi_{\theta}\left( a|s \right)}{\pi_{\theta_{old}}\left( a|s \right)}\)
is the probability ratio,
\({\widehat{A}}_{t} = G_{t} - \upsilon_{\pi}(s)\) is the estimated
advantage, and \(\epsilon\) is a small hyperparameter controlling the
clipping range. This method ensures that the updates are stable and
avoid excessively large policy changes (\colorbox{yellow!40}{Reference}). Gradient-free
methods, on the other hand, typically rely on search-based techniques or
evolutionary strategies to update the policy without directly computing
gradients (\colorbox{yellow!40}{Reference}). These methods are particularly useful in
environments where the gradient is difficult to compute or highly noisy
(\colorbox{yellow!40}{Reference}). More recently, both model-based and model-free methods
have incorporated deep learning, leading to algorithms like Actor-Critic
and Deep Q-Network (\colorbox{yellow!40}{e.g., Reference 1, Reference 2}), which have
paved the way for a new branch of RL known as deep RL (\colorbox{yellow!40}{cf.,
Reference}). The ability to make autonomous sequential decisions and
adapt to stochastic environments makes RL particularly well-suited for
production scheduling applications (\colorbox{yellow!40}{cf., Reference 1, Reference 2}),
providing a foundation for the remainder of this work.
